% Working draft of the transitivity lemma, curated in chat.
% MOVEMENT 1 of 4, revision 2: the Let, the two states, the tape and the disk
% action visible, and the block reorganized around its logical core --
% one element, invertible at N, hence orbit-stabilizer at every t, hence the
% sweep.  Not yet in the master.

\begin{proof}
Let
\[
  P,Q\in\int_{\mathcal B}\mathcal R_A
\]
be arbitrary. By the inverse-image definition, the fibre data of $P$ and $Q$
supply north residue states $x_N,y_N$ together with base arrows
\[
  g:N\longrightarrow P.\mathrm{base},\qquad
  h:N\longrightarrow Q.\mathrm{base},
\]
whose A-section transports recover the two given fibre states. So the arbitrary
problem reduces to a north endomorphism $k:N\to N$ and a fibre morphism
\[
  \varphi:\mathcal A_A(k)(x_N)\longrightarrow y_N .
\]
Once these are in hand, functorial transport along
$g^{-1}\fatsemi k\fatsemi h$ gives an ambient Grothendieck morphism $P\to Q$,
and the full inclusion $\int\iota_A$ returns that morphism to
$\int_{\mathcal B}\mathcal R_A$. It remains only to construct the north
comparison.

\medskip
\noindent\emph{The two states.}
Both are selected by the C3 zero-sphere system, so each carries a slice
direction and its positioned coordinate is an authored C3 coordinate:
\[
  x_N=(I_1,z_1),\qquad y_N=(I_2,z_2),
  \qquad I_1,I_2\in S^6 .
\]
The north object frame is the distinguished disk action $D_A$ itself, so each
positioned coordinate has exactly one input behind it,
\[
  z_1=D_A\mathbin{\cdot}u_1,\qquad z_2=D_A\mathbin{\cdot}u_2 .
                                    \tag{P}\label{eq:north-preimages}
\]
These are the two objects to be joined, and they differ in nothing but their
inputs.

\medskip
\noindent\emph{What built them.}
Each is one evaluation of the A-section action functor along the fixed tape.
The tape is the $0$-to-$N$ transport of the inverse-image construction, run once
for each state,
\[
  \operatorname{eulerToNorth}_1,\ \operatorname{eulerToNorth}_2
    \;:\;0\longrightarrow N,
\]
and at $N$ the functor records
\[
  (I,u)\longmapsto
  \bigl((I,u),\;D_A\mathbin{\cdot}(I,u),\;
    A_{\OO}(D_A\mathbin{\cdot}(I,u))\bigr).
                                    \tag{A}\label{eq:north-graph}
\]
Run at the two inputs, with the disk action left standing:
\[
\begin{array}{l}
  (I_1,u_1)\longmapsto
    \bigl((I_1,u_1),\;D_A\mathbin{\cdot}(I_1,u_1),\;
      A_{\OO}(D_A\mathbin{\cdot}(I_1,u_1))\bigr),\\[0.9ex]
  (I_2,u_2)\longmapsto
    \bigl((I_2,u_2),\;D_A\mathbin{\cdot}(I_2,u_2),\;
      A_{\OO}(D_A\mathbin{\cdot}(I_2,u_2))\bigr).
\end{array}
\]
The left entry is the input, the middle entry is the positioned coordinate, and
the right entry is the realized value. Reading the middle entries gives
\eqref{eq:north-preimages}, which is the residue condition; reading the left
entries gives the inputs the comparison is evaluated at.

\medskip
\noindent\emph{Why one element suffices.}
There is one construction here and not two. Euler presents at $0$, and the same
presentation arriving at $N$ \emph{is} Weierstrass: \textup{C1} continues $g_A$
through the simple pole and the Weierstrass divisor is exact at $N$, so at $N$
the two expressions are the same function. That is what makes the common datum
nonzero, $u_A=g_A(p_A)\neq0$ \textup{(}Definition~\ref{def:base}\textup{)}, and
hence what makes $D_A$ invertible: $\operatorname{diag}(u_A,1)$ has determinant
$u_A\neq0$, its class lies in the projective group, and $D_A$ is its conjugate
by the Cayley map. Invertibility at $N$ is not an assumption about $D_A$; it is
what the continuation delivers at the point where the two presentations agree.

Because $D_A$ is invertible, every object frame
$a_A(b)=\operatorname{cayleyProjective}(o_b)D_A$ is a group element, and
\eqref{eq:orbit-stabilizer} applies at every object of the base --- in
particular at every instant of the tape, up to and including $N$. The one
element positioned everywhere is what the sweep sweeps.

\medskip
\noindent\emph{What the sweep gives.}
The direction $I$ ranges over $S^6$ and $G_2$ carries it; $D_A$ positions the
coordinate $u$ in the slice that direction selects; and $A_{\OO}$ realizes the
value there. One element, swept over all directions, is the whole content of
$\mathcal A_A$, and its well-definedness is \eqref{eq:orbit-stabilizer} and
nothing else. The inverse image is taken over exactly what is swept, which is
why the comparison below is available in $\int_{\mathcal B}\mathcal R_A$ and not
in the ambient total.
\end{proof}
